% STATUS: written 2026-11-06 (Ferrero-Haetinger 2002, Narvaez Macarro 2009,
% Hasse-Schmidt 1937 and Nakai 1970 read in full)
% Additions of 2026-11-07 moved from chapter 20: Narvaez Macarro 2020
% (Introduction, Definition 2.1.1 p. 15, pp. 1--2) read in full.
% STATUS: extended 2026-11-06 (Hosseini and Rehman 2019, Nagamine 2019 read in full from the library PDFs)
% STATUS: extended 2026-11-07 (Tirado Hernandez 2021 read in full from the
% library PDF; Barajas and Duarte 2026 read in full while adding chapter 11's
% trivial-extension section)
% Sources read in full: Ferrero--Haetinger 2002 (pp. 2321--2333); Narvaez Macarro 2009 (pp. 2979--3014, section 1.1);
% Hasse--Schmidt 1937, J. Reine Angew. Math. 177, pp. 215--237 (the Zusatz bei der
% Korrektur, pp. 223--237, holds the ring-theoretic definition, p. 224, and Satze 1
% and 2, p. 226); Nakai 1970, Osaka J. Math. 7, pp. 1--27 (definition p. 1,
% Theorem 2 p. 3, Proposition 5 p. 4).
% Tirado Hernandez 2021: the PDF has no printed journal page numbers, so
% citations give the page of the PDF. Read 2026-11-07: the definition of
% integrability (p. 1) and Proposition 2.1 (p. 4).
% Ashraf and Ansari 2021, Moafian 2017 and Ashraf and Jabeen 2018 read in full
% 2026-11-07: the definitions of a Lie higher derivation, the properness
% question, and the splitting theorems for triangular and generalized matrix
% algebras.
% Origins: the definition is quoted from Ferrero--Haetinger 2002 (p. 2323, after Jacobson,
% Basic Algebra II, Exercise 4, p. 532), the two rules of the sequence are recalled
% from Hasse--Schmidt 1937 first-hand (p. 224), and the single-map notion from
% Nakai 1970 first-hand (p. 1). Nakai attributes that notion to Osborn, Math. Ann.
% 170 (1967), 221--244; Osborn's paper is not in this book's collection, so it is
% named here through Nakai's attribution and not quoted.
\section{Higher derivations and Hasse--Schmidt sequences}\label{sec:higher-multi:hasse-schmidt}

Every generalization in Chapter~\ref{ch:generalized-twisted} changed one
map. The generalization of this chapter replaces one map by a sequence
\cite[p.~2323]{FerreroHaetinger2002}. The first component of the sequence is
the identity, and the next is a derivation
\cite[p.~2323]{FerreroHaetinger2002}. The product rule of the whole sequence
collects
the terms $D_i(x)D_j(y)$ along the \emph{antidiagonal} $i+j=n$, the terms
whose two indices sum to $n$ \cite[p.~2323]{FerreroHaetinger2002}.

\begin{definition}[{\cite[p.~2323]{FerreroHaetinger2002}}]\label{def:higher:hd}
Let $T$ be a ring. A family of additive mappings
$D=(d_i)_{i\in\mathbb{N}}$ of $T$ is a \emph{higher derivation} of $T$ if
$d_0=\mathrm{id}_T$ and, for every $n\in\mathbb{N}$,
\begin{equation*}
d_n(ab)=\sum_{i=0}^{n} d_i(a)\,d_{n-i}(b)
\qquad\text{for all }a,b\in T.
\end{equation*}
\end{definition}

The definition reaches the literature as an exercise
\cite[p.~2323]{FerreroHaetinger2002}. Ferrero and Haetinger record it with
page 532, Exercise 4 of the book \emph{Basic Algebra~II} by Jacobson
\cite[p.~2323]{FerreroHaetinger2002}.

Hasse and Schmidt gave the same sequence for algebraic function fields of one
variable in 1937 \cite[pp.~215--223]{HasseSchmidt1937}. Their paper states the
two rules of the sequence for an arbitrary integral domain
\cite[pp.~223--224]{HasseSchmidt1937}. The authors set a domain $I$ of
characteristic $p$ and a larger domain $I^{*}$ containing the derivatives
\cite[p.~224]{HasseSchmidt1937}. They assign to each $y\in I$ a function
$y^{(\nu)}\in I^{*}$ for every natural number $\nu$
\cite[p.~224]{HasseSchmidt1937}. The assignment is a \emph{differentiation}
when it satisfies the two rules below \cite[p.~224]{HasseSchmidt1937}.

\begin{definition}[{\cite[p.~224]{HasseSchmidt1937}}]\label{def:higher:hsrules}
Let $I\subseteq I^{*}$ be integral domains of characteristic $p$, and let
$y\mapsto y^{(\nu)}$ assign to each $y\in I$ and each $\nu\in\mathbb{N}$ an
element $y^{(\nu)}\in I^{*}$. The assignment is a \emph{differentiation} of
$I$ when it satisfies the sum rule
\[
(y+z)^{(\nu)}=y^{(\nu)}+z^{(\nu)}
\]
and the modified product rule
\[
(yz)^{(\nu)}=y^{(\nu)}z+y^{(\nu-1)}z'+\cdots+y'z^{(\nu-1)}+yz^{(\nu)}
\]
for all $y,z\in I$ and all $\nu$. The element $y^{(\nu)}$ is the
\emph{$\nu$-th derivative} of $y$.
\end{definition}

A function whose derivatives all vanish from some order on is a
\emph{differentiation constant} of that order, and the absolute constants are
those whose derivatives all vanish \cite[p.~224]{HasseSchmidt1937}. The authors
note that the constants of at least a given order form a subdomain of $I$, and
that when $I$ is a field they form a subfield
\cite[p.~224]{HasseSchmidt1937}.

The 1937 paper then identifies a sequence with a single map into a ring of
power series \cite[p.~225]{HasseSchmidt1937}. Over $I^{*}$ it forms the domain
$T^{*}=I^{*}\{u\}$ of all power series in $u$, and assigns to each $y\in I$
the series $Y_{u}=y+y'u+y''u^{2}+\cdots$ \cite[p.~225]{HasseSchmidt1937}. The
authors call $Y_{u}$ the general Taylor expansion of $y$
\cite[p.~225]{HasseSchmidt1937}. They prove that the assigned series form a
subdomain on which the assignment is an isomorphism
\cite[Satz~1, p.~226]{HasseSchmidt1937}. They prove the converse too
\cite[Satz~2, p.~226]{HasseSchmidt1937}. Take an isomorphism of $I$ onto a
subdomain of $T^{*}$ that carries each $y$ to a series with first term $y$.
Reading the coefficients as derivatives gives a differentiation
\cite[Satz~2, p.~226]{HasseSchmidt1937}.

Nakai took up the single map rather than the sequence in 1970
\cite[p.~1]{Nakai1970}. He works over commutative rings $k$ and $A$ with $A$ a
$k$-algebra, and he attributes the notion to Osborn
\cite[p.~1]{Nakai1970}. His $q$-th order derivation keeps the product rule on
$(q+1)$ arguments, and his Theorem~2 turns that rule into an identity of
commutators \cite[Theorem~2, p.~3]{Nakai1970}. The link between the two
notions is his Proposition~5: the $q$-th component of a higher derivation is a
$q$-th order derivation \cite[Proposition~5, p.~4]{Nakai1970}. A $q$-th order
derivation need not be a component of a higher derivation, and Nakai leaves
that question open \cite[p.~1]{Nakai1970}.

The model is the iterated derivative of a single derivation
\cite[p.~2322]{FerreroHaetinger2002}. If $d$ is a derivation of a rational
algebra, the sequence $d_i=d^{i}/i!$ is a higher derivation of $R$
\cite[p.~2322]{FerreroHaetinger2002}.

The name Hasse--Schmidt marks a second presentation of the same notion
\cite[p.~2983]{NarvaezMacarro2009}. Narváez Macarro states it for algebras,
with the 1937 paper as his reference for the definition
\cite[§1.1, reference~7]{NarvaezMacarro2009}.

\begin{definition}[{\cite[§1.1, p.~2983]{NarvaezMacarro2009}}]\label{def:higher:hs}
Let $k$ be a commutative ring, $f\colon k\to A$ and $g\colon A\to B$ ring
homomorphisms, and $m\geq 1$ an integer or $m=\infty$. A
\emph{Hasse--Schmidt derivation} over $k$ of length $m$ from $A$ to $B$ is a
sequence $D=(D_0,D_1,\dots,D_m)$ or $D=(D_0,D_1,\dots)$ of $k$-linear
maps $D_i\colon A\to B$ with
\begin{equation*}
D_0=g,\qquad D_i(xy)=\sum_{r+s=i}D_r(x)\,D_s(y)
\qquad\text{for all }x,y\in A\text{ and all }i.
\end{equation*}
When $A=B$ and $g=\mathrm{id}_A$ one says that $D$ is a Hasse--Schmidt
derivation of $A$ over $k$.
\end{definition}

Such a sequence is the same data as one algebra map into a truncated
polynomial ring \cite[p.~2984]{NarvaezMacarro2009}. With
$B_m=B[t]/(t^{m+1})$, the sequence $D$ is determined by the $k$-algebra
homomorphism $\Phi\colon A\to B_m$ given by
$\Phi(x)=\sum_{i=0}^{m} D_i(x)\,t^i$ \cite[p.~2984]{NarvaezMacarro2009}.

Over a semiprime ring a higher derivation extends to the Martindale ring of
quotients in a unique way \cite[p.~2323]{FerreroHaetinger2002}.

\begin{proposition}[{\cite[Proposition~1.1, p.~2323]{FerreroHaetinger2002}}]\label{prop:higher:extend}
Let $R$ be a semiprime ring and $D=(d_i)_{i\in\mathbb{N}}$ a higher
derivation of $R$. Then there exists a unique higher derivation
$D^{*}=(d^{*}_i)_{i\in\mathbb{N}}$ of $Q$ such that
$d^{*}_n|_R=d_n$ for every $n\in\mathbb{N}$.
\end{proposition}

The theory of the sequence starts when the terms satisfy a linear relation
\cite[p.~2323]{FerreroHaetinger2002}. A family $D$ satisfies an $S$-linear
relation on $T$ when there exist $a_0,a_1,\dots,a_n\neq 0$ in $S$ with
$\sum_{i=0}^{n} a_i d_i(x)=0$ for all $x\in T$; the minimal $n$ is the
length of $D$. For a single derivation $d$ of a rational algebra the
relation is the algebraicity of $d$, that is, the existence of a polynomial
relation among its powers \cite[p.~2322]{FerreroHaetinger2002}. On a prime
ring the relation is rigid, as Corollary~\ref{cor:higher:prime} shows
\cite[p.~2325]{FerreroHaetinger2002}.

\begin{corollary}[{\cite[Corollary~1.4, p.~2325]{FerreroHaetinger2002}}]\label{cor:higher:prime}
Assume that $R$ is a prime ring and $D=(d_i)_{i\in\mathbb{N}}$ a higher
derivation of $R$ that satisfies an $R$-linear relation of minimal length
$n$ on $R$. Then there exist $q_0=1,q_1,\dots,q_{n-1}\in Q$ such that
$\sum_{i=1}^{n} q_{n-i} d_i=0$. In particular, $D$ satisfies a monic
$Q$-linear relation of length $n$ on $R$. Moreover, $d_1=d_{q_1}$ and
\begin{equation*}
d_m(x)=d_{q_m}(x)-\sum_{i=1}^{m-1} q_i\,d_{m-i}(x)
\qquad\text{for all }x\in R\text{ and }2\leq m\leq n.
\end{equation*}
\end{corollary}

The main result compares the possible relations
\cite[pp.~2327--2328]{FerreroHaetinger2002}. It says the ring of
coefficients of the relation does not matter, as
Theorem~\ref{thm:higher:relations} states
\cite[p.~2327]{FerreroHaetinger2002}.

\begin{theorem}[{\cite[Theorem~1.8, pp.~2327--2328]{FerreroHaetinger2002}}]\label{thm:higher:relations}
Assume that $R$ is a semiprime ring and $D$ a higher derivation of $R$.
Then the following conditions are equivalent:
\begin{enumerate}[label=(\roman*)]
\item $D$ satisfies an $R$-linear relation on $R$;
\item $D$ satisfies a $Q$-linear relation on $R$;
\item $D^{*}$ satisfies a $Q$-linear relation on $Q$;
\item $D^{*}$ satisfies an $R$-linear relation on $Q$.
\end{enumerate}
Moreover, when one of these conditions holds, the minimal lengths of the
relations are all the same.
\end{theorem}

A higher derivation is therefore a derivation whose values are measured in
steps, as Definition~\ref{def:higher:hd} displays. Later sections of this
chapter multiply the map in a second direction, two derivation arguments at
once in Section~\ref{sec:higher-multi:biderivations} and three in
Section~\ref{sec:higher-multi:tri-generalized}.

\section{Coding theory: Hasse--Schmidt derivations in positive characteristic}\label{sec:higher-multi:coding}

Error-correcting codes are built over finite fields, where the classical
derivative no longer records enough information, and Hasse--Schmidt
derivations replace it \cite[Introduction]{NarvaezMacarro2020}. Narv\'aez
Macarro states the multivariate form
\cite[Definition~2.1.1, p.~15]{NarvaezMacarro2020}.

\begin{definition}[{\cite[Definition~2.1.1, p.~15]{NarvaezMacarro2020}}]\label{def:higher:hasse-schmidt-variate}
Let $k$ be a commutative ring, $A$ a commutative $k$-algebra, $\mathbf{s}=\{s_1,\dots,s_p\}$
a set of variables and $\Delta\subset\mathbb{N}^p$ a non-empty co-ideal. A
\emph{$(p,\Delta)$-variate Hasse--Schmidt derivation} of $A$ over $k$ is a
family $D=(D_\alpha)_{\alpha\in\Delta}$ of $k$-linear maps $D_\alpha\colon A\to A$
satisfying
\[
D_0=\mathrm{Id}_A,\qquad D_\alpha(xy)=\sum_{\beta+\gamma=\alpha}D_\beta(x)D_\gamma(y)
\quad\text{for every }x,y\in A\text{ and every }\alpha\in\Delta.
\]
For $p=1$ and $\Delta=\{q\in\mathbb{N}\mid q\leq m\}$, $D$ is a
\emph{Hasse--Schmidt derivation of length $m$}, also called a higher derivation.
\end{definition}

The first component $D_1$ is a $k$-derivation of $A$ in the sense of
Chapter~\ref{ch:basic-theory}. Every $D_i$ for $i\geq 1$ is a $k$-linear
differential operator of order $\leq i$ vanishing on $1$
\cite[p.~1]{NarvaezMacarro2020}. In characteristic $0$, every $k$-derivation
$\delta$ of $A$ appears as the $1$-component of a Hasse--Schmidt derivation of
infinite length, namely the exponential $e^{t\delta}=\sum_{\alpha\geq0}
\frac{\delta^\alpha}{\alpha!}t^\alpha$ \cite[p.~2]{NarvaezMacarro2020}. In
positive characteristic this exponential collapses when the derivation reaches
order $p$, and the higher derivations are the surviving notion \cite[p.~2]{NarvaezMacarro2020}.

\hbfigure{07-higher-multi/hasse-schmidt-collapse}{Why the exponential $e^{t\delta}=\sum_{\alpha\ge0}\delta^{\alpha}t^{\alpha}/\alpha!$ survives in characteristic $0$ and collapses in characteristic $p$ \cite[p.~2]{NarvaezMacarro2020}. The coefficients $1/\alpha!$ exist for every $\alpha$ in characteristic $0$. In characteristic $p$ the factorial $p!$ is a multiple of $p$, so $1/p!$ does not exist and the series stops at $\alpha=p$; the higher derivations are the notion that survives.}{fig:higher:collapse}
The rule can be weakened by a pair of maps, one for each side, and the sequence
keeps its shape \cite[Definition~2.1, p.~3]{HosseiniRehman2019Higher}. Hosseini
and Rehman call a sequence a higher
$\{g_n,h_n\}$-derivation when each term splits the product according to two
sequences $\{g_n\}$ and $\{h_n\}$ \cite[Definition~2.1, p.~3]{HosseiniRehman2019Higher}.
They express every such sequence by a sequence of $\{G,H\}$-derivations
\cite[Theorem~2.4, p.~7]{HosseiniRehman2019Higher}. A Jordan higher
$\{g_n,h_n\}$-derivation on a semiprime algebra is then a higher
$\{g_n,h_n\}$-derivation \cite[Corollary~2.6, p.~12]{HosseiniRehman2019Higher}.
The result extends the semiprime theorem of Brešar for the one-term
$\{g,h\}$-derivation \cite[p.~1]{HosseiniRehman2019Higher}.

The same sequence has a Jacobian form in positive characteristic
\cite[p.~1]{Nagamine2019HigherJacobian}. Nagamine defines the higher derivation of
Jacobian type by a tuple $F=(f_1,\dots,f_{n-1})$ of polynomials
\cite[Definition~2.4, p.~3]{Nagamine2019HigherJacobian}. The tuple gives a sequence
of $R$-linear maps, and the sequence is a higher derivation exactly when the
tuple extends to a coordinate system
\cite[Definition~2.4, p.~3]{Nagamine2019HigherJacobian}. An $(n-1)$-tuple of
polynomials is extendable in this sense exactly when the higher derivation
exists and has a slice \cite[Theorem~2.8, p.~5]{Nagamine2019HigherJacobian}. For the
polynomial ring in two variables over a field of positive characteristic, the
same construction describes the variables \cite[Theorem~3.1, p.~6]{Nagamine2019HigherJacobian}.

A single derivation need not be the first component of a sequence, and the
question of which ones are is the question of integrability
\cite[p.~1]{Tirado2021}. A $k$-derivation $\delta$ of $A$ is \emph{$m$-integrable}
when there is a Hasse--Schmidt derivation $D$ of length $m$ with $D_{1}=\delta$
\cite[p.~1]{Tirado2021}. The set of $m$-integrable derivations is a module inside
the module of derivations, written $\mathrm{IDer}_{k}(A;m)$
\cite[p.~1]{Tirado2021}. In characteristic zero, or when $A$ is smooth over $k$,
every derivation is integrable \cite[p.~1]{Tirado2021}. In positive characteristic
these modules carry more information than the module of all derivations
\cite[p.~1]{Tirado2021}. Tirado Hern\'andez computes them for the plane curve
$x^{n}=y^{q}$, in each of several cases of the exponents
\cite[Proposition~2.1, p.~4]{Tirado2021}.

Barajas and Duarte tie the sequence to the trivial extension of a ring by a module
\cite[p.~2]{BarajasDuarte2026}. A derivation from $A$ into $M$ is the same data as a
ring homomorphism from $A$ into the trivial extension $A\ltimes M$ that the
projection takes back to the identity \cite[p.~2]{BarajasDuarte2026}. Anderson,
Bennis, Fahid and Shaiea had replaced the one module by $n$ modules
\cite[p.~2]{BarajasDuarte2026}. Barajas and Duarte replace the derivation by a
sequence of maps with the same property, and call it a multi Hasse--Schmidt
derivation \cite[Definition~3.1, p.~6]{BarajasDuarte2026}.

\begin{proposition}[{\cite[Proposition~3.5, p.~7]{BarajasDuarte2026}}]\label{prop:higher:multi-hs}
Let $A,B_{1},\dots,B_{n}$ be $k$-algebras, let $g_{i}\colon A\to B_{i}$ be $k$-algebra
homomorphisms, and give each $B_{i}$ the $A$-module structure induced by $g_{i}$.
Then the multi Hasse--Schmidt derivations of $A$ in $B_{1},\dots,B_{n}$ correspond
bijectively to the $k$-algebra homomorphisms from $A$ into the $n$-trivial
extension $A\ltimes_{n}B_{1}\ltimes\cdots\ltimes B_{n}$ whose components all agree
with the $g_{i}$.
\end{proposition}

When all the algebras $B_{i}$ are the same, the multi Hasse--Schmidt derivations
are the Hasse--Schmidt derivations of Definition~\ref{def:higher:hs}
\cite[p.~2]{BarajasDuarte2026}. For $n=1$ the correspondence is the classical one
between derivations and trivial extensions \cite[p.~2]{BarajasDuarte2026}.

The sequence has a Lie version. A Lie higher derivation of an algebra $A$ is a
sequence $L=(L_{n})$ of linear maps with $L_{0}$ the identity and
$L_{n}([x,y])=\sum_{i+j=n}[L_{i}(x),L_{j}(y)]$ for all $x,y$
\cite[p.~1]{AshrafAnsari2021}. The first term is then a Lie derivation
\cite[p.~1]{AshrafAnsari2021}. The authors ask when such a sequence is proper, that
is, when it is a higher derivation plus a sequence of central maps that vanish on
commutators \cite[p.~1]{AshrafAnsari2021}.

\begin{theorem}[{\cite[Theorem~3.1, p.~6]{AshrafAnsari2021}}]\label{thm:higher:lie-higher-triangular}
Let $A=\mathrm{Tri}(A,M,B)$ be a triangular algebra and $L=(L_{n})$ a Lie higher
derivation of $A$. Then there is a triangular algebra $A_{0}$ containing $A$ with
the same identity such that $L_{n}=\Delta_{n}+\chi_{n}$ for every $n$, where
$(\Delta_{n})$ is a higher derivation of $A$ into $A_{0}$ and each $\chi_{n}$ is a
linear map of $A$ into the center of $A_{0}$ that vanishes on every commutator.
\end{theorem}

Moafian proved the corresponding structure theorem for a generalized matrix
algebra, and he states conditions under which the sequence is proper
\cite[Theorem~4.4, p.~12]{Moafian2017}. Ashraf and Jabeen carry the same splitting
to a nonlinear generalized Lie triple higher derivation
\cite[p.~1]{AshrafJabeen2018}. There each component is the sum of an additive
generalized higher derivation and a nonlinear map that vanishes on all Lie triple
products \cite[Theorem~3.1, p.~5]{AshrafJabeen2018}.

Over a domain of positive characteristic the higher derivations of Jacobian
type carry the same information as the Jacobian matrix. Nagamine gives a
necessary and sufficient condition for an $(n-1)$-tuple of such derivations on
$R[x_1,\dots,x_n]$ \cite[Abstract, p.~1]{Nagamine2019HigherJacobian}.

Hasse-Schmidt derivations are put in the language of the Hopf algebra of non-commutative symmetric functions, which is the algebraic form the chapter recalls from Hasse and Schmidt~\cite{WlHasseSchmidtDerivations2012}.

Orthogonal generalized higher derivations on a semiprime ring are studied, the orthogonality being imposed on the sequences~\cite{WlOrthogonalGeneralizedHigherDer2021}.

The kernels of locally finite higher derivations are studied, which is the constants object for a sequence rather than for one derivation~\cite{WlNotesOnTheKernelsOfLocallyFini2018}.

For a skew higher derivation $(d_i)$ on a semiprime ring with a suitable torsion restriction, Rajapati proves that every $n$-centralizing skew higher derivation is $n$-commuting, and that if a ring carries an $n$-centralizing skew higher derivation then either the ring is commutative or some linear combination of the $d_i$ maps the centre to zero~\cite{WlNCommutingSkewHigherDerivation2020}.

Das calls an algebra together with a distinguished higher derivation an AssHDer pair, defines a cohomology for these pairs with coefficients in a representation, and shows that isomorphism classes of central extensions are classified by the second cohomology group with trivial coefficients and those of abelian extensions by the second cohomology with an arbitrary representation~\cite{WlExtensionsAndDeformationsOfAlg2022}.

Fei and Zhang study nonlinear higher anti-derivable mappings on generalized matrix algebras and show that they are additive under stated conditions~\cite{WlAdditivityOfNonlinearHigherAnt2023}.

Khazil proves results on ideals of semiprime rings carrying orthogonal higher derivations, the main stated conclusion being that for higher derivations $f_n$ and $h_n$ the product $f_nh_n$ vanishes~\cite{WlIdealOfSemiPrimeRingsWithOrtho2021}.
